\Unit[Reflection; refraction (Snell); thin lens and magnification; interference/diffraction]{Optics}
\label{unit:optics}

\begin{KnowledgePoint}{Reflection and Refraction \Core}
  \KPFormulas
  \begin{formulabox}
  \textbf{Concept explanation:} Light reflects with equal incident and reflected angles; refraction across media obeys Snell's law.

  \textbf{Core formulas:}
  \[
  \left\{\begin{aligned}
    &\theta_i=\theta_r,\\
    &n_1\sin\theta_1=n_2\sin\theta_2,\quad \text{TIR when }\theta_1>\theta_c=\arcsin(n_2/n_1)\ (n_1>n_2).
  \end{aligned}\right.
  \]

  \textbf{Variable definitions:} $n$ refractive index; $\theta$ angles measured to the normal.

  \textbf{Prerequisites \& scope:} Geometric optics regime; isotropic media; polarization effects ignored here.
  \end{formulabox}

  \KPHeuristics
  \begin{heuristicsbox}
  \begin{itemize}[leftmargin=*]
    \item Draw the normal and principal rays first; search for total internal reflection when going to a lower-$n$ medium.
    \item Use reversibility of light to validate constructions.
  \end{itemize}
  \end{heuristicsbox}

  \KPProblems
\begin{cheatproblem}
  Light travels from air ($n_1=1$) into water ($n_2=1.33$) at incidence angle $\theta_1=40^\circ$. Find the refraction angle $\theta_2$.
\begin{solutionbox}
  Snell: $n_1\sin\theta_1=n_2\sin\theta_2\Rightarrow\sin\theta_2=\dfrac{\sin40^\circ}{1.33}\Rightarrow\theta_2\approx28.9^\circ$.
\end{solutionbox}
\end{cheatproblem}
\begin{cheatproblem}
  Light moves from glass ($n=1.5$) to air ($n=1$). Find the critical angle for total internal reflection.
\begin{solutionbox}
  $\theta_c=\arcsin\dfrac{n_2}{n_1}=\arcsin\dfrac{1}{1.5}\approx41.8^\circ$.
\end{solutionbox}
\end{cheatproblem}
\end{KnowledgePoint}

\begin{KnowledgePoint}{Thin Lenses and Sign Conventions \Core}
  \KPFormulas
  \begin{formulabox}
  \textbf{Concept explanation:} Thin lens imaging follows the lens equation with sign conventions; magnification uses image and object sizes/orientations.
  
  \textbf{Core formulas:}
  \[
  \left\{\begin{aligned}
    &\frac{1}{f}=\frac{1}{s}+\frac{1}{s'},\quad m=-\frac{s'}{s}=\frac{h'}{h}.
  \end{aligned}\right.
  \]

  \textbf{Variable definitions:} $f$ focal length; $s$ object distance; $s'$ image distance; $m$ magnification; $h',h$ image/object heights.

  \textbf{Prerequisites \& scope:} Use consistent sign convention (e.g., real is positive); paraxial approximation.
  \end{formulabox}

  \KPHeuristics
  \begin{heuristicsbox}
  \begin{itemize}[leftmargin=*]
    \item Combine equation + ray diagram: draw two principal rays to confirm the algebraic image location.
    \item Remember that negative $m$ indicates inversion; $|m|>1$ indicates magnification.
    \item Sign convention (real-is-positive): take $s>0$ for real objects and $s'>0$ for real images on the opposite side of the lens from the object; $s'<0$ indicates a virtual image on the object side (then $m>0$ and the image is upright).
  \end{itemize}
  \end{heuristicsbox}

  \KPProblems
  % Simple principal-ray diagram (converging lens, object inside f -> virtual image)
  \begin{center}
    \begin{tikzpicture}[x=0.06cm,y=0.06cm,>=Latex]
      % Optical axis
      \draw[-] (-60,0) -- (80,0);
      % Lens at x=0
      \draw[thick] (0,-25) -- (0,25);
      % Focal points at ±60 cm (f=60 cm)
      \fill (-60,0) circle (1.2pt) node[below=2pt] {$F$};
      \fill (60,0) circle (1.2pt) node[below=2pt] {$F'$};
      % Object (arrow) at s=30 cm left of lens (inside focal length)
      \draw[thick,->] (-30,0) -- (-30,18) node[left] {object};
      % Principal ray 1: parallel to axis, refracted as if through F'
      \draw[->] (-30,18) -- (0,18);      % incident
      \draw[->] (0,18) -- (60,0);         % refracted
      % Principal ray 2: through center (undeviated)
      \draw[->] (-30,18) -- (0,0);
      \draw[->] (0,0) -- (30,-18);
      % Back-projections of refracted rays to locate virtual image at s'=-60 cm
      \draw[dashed] (0,18) -- (-60,36);
      \draw[dashed] (0,0) -- (-60,36);
      % Image arrow (virtual, upright, magnified)
      \draw[thick,->,blue] (-60,0) -- (-60,36) node[left] {virtual image};
    \end{tikzpicture}
  \end{center}
\begin{cheatproblem}
  An object at $s=30\,\text{cm}$ forms an image at $s'=-60\,\text{cm}$ using a thin lens. Find the focal length $f$ and magnification $m$.
  \begin{solutionbox}
  Lens equation: $\tfrac{1}{f}=\tfrac{1}{s}+\tfrac{1}{s'}=\tfrac{1}{30}+\tfrac{1}{-60}=\tfrac{1}{60}$, so $f=60\,\text{cm}$.
  Magnification: $m=-s'/s = -(-60)/30=2$.
  Since $f>0$, it is a converging lens. Since $s'<0$, the image is virtual. Since $m>0$, the image is upright.
  \end{solutionbox}
  \end{cheatproblem}
\end{KnowledgePoint}

\begin{KnowledgePoint}{Interference and Diffraction \Core}
  \KPFormulas
  \begin{formulabox}
  \textbf{Concept explanation:} Coherent sources produce interference patterns; finite apertures make the light diffract, establishing angular scales by wavelength/aperture.

  \textbf{Core formulas:}
  \[
  \left\{\begin{aligned}
    &\text{Double-slit maxima: } d\sin\theta=m\lambda,\\
    &\text{Single-slit minima: } a\sin\theta=m\lambda,\ m=\pm1,\pm2,\dots
  \end{aligned}\right.
  \]

  \textbf{Variable definitions:} $d$ slit separation; $a$ slit width; $\lambda$ wavelength; $\theta$ diffraction angle.

  \textbf{Prerequisites \& scope:} Small-angle approximations $\sin\theta\approx\theta$ valid near axis; coherence required for stable fringes.
  \end{formulabox}

  \KPHeuristics
  \begin{heuristicsbox}
  \begin{itemize}[leftmargin=*]
    \item Map angles to screen positions with $y\approx L\tan\theta\approx L\theta$ for small $\theta$.
    \item To resolve features, compare $\lambda$ to $a$ and $d$ to predict fringe spacing/envelope width.
  \end{itemize}
  \end{heuristicsbox}

  \KPProblems
  \begin{cheatproblem}
  For double-slit with spacing $d$ and wavelength $\lambda$, what is the angle of the $m$-th bright fringe?
  \begin{solutionbox}
  $d\sin\theta=m\lambda \Rightarrow \theta=\arcsin(m\lambda/d)$ (small-angle: $\theta\approx m\lambda/d$).
  \end{solutionbox}
  \end{cheatproblem}
\end{KnowledgePoint}


\UnitProblemLinks{optics}
